\documentclass[12pt,journal,compsoc]{IEEEtran}

%-----PACKAGES-------------------------------------------------------------------------------------

\usepackage{graphicx}        % For including graphics
\usepackage{amsmath}         % For mathematical formulas
\usepackage{cite}            % For citations
\usepackage{hyperref}        % For hyperlinks
\usepackage{array}           % For better arrays (e.g., matrices) in maths
\usepackage{caption}         % For customizing captions
\usepackage{subcaption}      % For subfigures
\usepackage{booktabs}        % For better table formatting
\usepackage{multicol}        % For multi-column text
\usepackage{enumitem}        % For customizable list environments
\usepackage{float}           % For improved float control
\usepackage{lipsum}          % For placeholder text

%----- The DOCUMENT Environment-------------------------------------------------------------------

\begin{document}

\title{Literature Review on Adaptive Telemetry Systems for Monitoring Neonatal Neurodevelopment}
\author{Shanaya I. Malik}


\date{\today}

% The paper headers
\markboth{Literature Review on Adaptive Telemetry Systems for Monitoring Neonatal Neurodevelopment}%

\IEEEcompsoctitleabstractindextext{%
\begin{abstract}
%\boldmath
The abstract goes here.
\end{abstract}

\begin{IEEEkeywords}
Adaptive, telemetry, systems, neonatal, neurodevelopment. 
\end{IEEEkeywords}}

\maketitle

%----- The SECTION Environment -------------------------------------------------------------------
\section{Introduction}

\IEEEPARstart{T}{his} demo file is intended to serve as a ``starter file''
for IEEE Computer Society journal papers produced under \LaTeX\ using
IEEEtran.cls version 1.7 and later.
% You must have at least 2 lines in the paragraph with the drop letter
% (should never be an issue)
Be sure to omit this text and insert your own Introduction text for your tutorial.

\subsection{Subsection Heading Here}
Subsection text here.

% Creating a subsubsection:

\subsubsection{Subsubsection Heading Here}
Subsubsection text here.


%----- Additional Features -----------------------------------------------------------------------

\section{Additional Features}

\subsection{Figures}
% FIGURES:

\begin{figure}[h] 	% There are several different modifiers that can be used in [].
\centering
\includegraphics[width=1.5in]{slug.pdf}
\caption{Sammy the Slug}
\label{fig_slug}
\end{figure}

%-------------------------------------------------------------------------------------------------
\subsection{Label, Cite, and Ref Commands}
You will need to be able to define and explain how to use the following commands:
\begin{verbatim} \label{fig_slug} \end{verbatim} as it appears in the figure environment\\
\begin{verbatim} \cite{IEEEhowto:kopka} \end{verbatim} appears like: \cite{IEEEhowto:kopka}\\
\begin{verbatim} \ref{fig_slug} \end{verbatim} appears like: \ref{fig_slug}\\ 

%-------------------------------------------------------------------------------------------------
\subsection{Tables}
An example of a floating table.
%
\begin{table}[h]
%% increase table row spacing, adjust to taste
\renewcommand{\arraystretch}{1.3}
% if using array.sty, it might be a good idea to tweak the value of
%\extrarowheight as needed to properly center the text within the cells
\caption{An Example of a Table}
\label{table_example}
\centering
%% Some packages, such as MDW tools, offer better commands for making tables
%% than the plain LaTeX2e tabular which is used here.
\begin{tabular}{|c||c|}
\hline
One & Two\\
\hline
Three & Four\\
\hline
\end{tabular}
\end{table}


\section{Conclusion}
Conclusion goes here.

%----- APPENDICES --------------------------------------------------------------------------------
\appendices
\section{Appendix Title}
Appendix one text goes here.

% you can choose not to have a title for an appendix
% if you want by leaving the argument blank
\section{}
Appendix two text goes here.


%----- ACKNOWLEDGEMENT SECTION -------------------------------------------------------------------
% Explain what the asterisk * does in the next line: 
\section*{Acknowledgements}

The author would like to thank...\\ \\

\begin{thebibliography}{1}

\bibitem{IEEEhowto:toet}
M.~C.~Toet and P.~M.~A.~Lemmers, “Brain Monitoring in Neonates,” \emph{Web of Science}, accessed May 14, 2024. [Online]. Available: \url{https://www.webofscience.com/wos/woscc/full-record/WOS:000263512700003}

\bibitem{IEEEhowto:nicklin}
S.~Nicklin, A.~et al., “Neonatal Intensive Care Monitoring,” \emph{Current Paediatrics}, accessed May 14, 2024. [Online]. Available: \url{https://www.sciencedirect.com/science/article/pii/S0957583903001283}

\bibitem{IEEEhowto:poveda}
J.~Poveda, M.~O'Sullivan, E.~Popovici, and A.~Temko, "Portable neonatal EEG monitoring and sonification on an Android device," in \emph{Proc. 2017 39th Annual International Conference of the IEEE Engineering in Medicine and Biology Society (EMBC)}, Jeju, Korea (South), 2017, pp. 2018-2021, doi: 10.1109/EMBC.2017.8037248.

\end{thebibliography}

% if you will not have a photo at all:
\begin{IEEEbiographynophoto}{John Doe}
Biography text here.
\end{IEEEbiographynophoto}

% insert where needed to balance the two columns on the last page with
% biographies
%\newpage

\begin{IEEEbiographynophoto}{Jane Doe}
Biography text here.
\end{IEEEbiographynophoto}

\end{document}
